\documentclass[11pt,a4paper,notitlepage]{article}

%usepackage inclusions:
\usepackage{mathtools} %package that allows math formatting
\usepackage{commath} %package that allows for nice differential operators
\usepackage{bigints} %package that allows for bigger integral signs
\usepackage{amsfonts}   %package that allows for more math characters
\usepackage{amssymb}  %package that allows for more math symbols
\usepackage{shadethm}  %package that allows for shaded definitions&proofs
\usepackage[usenames,dvipsnames,svgnames,table]{xcolor} %package that allows colored text
\usepackage{booktabs}
\usepackage{float} % by Jakob
\usepackage[table]{xcolor} %to color table rows & columns
\usepackage{setspace} %to allow for the insertion line spaces at will
\usepackage{wrapfig} %to allow for figures to be captioned
\usepackage{mdframed} % to allow for wide frames and boxes around math&text
\usepackage[makeroom]{cancel} %to allow for strike-out text
\usepackage{subfigure}
\usepackage{parskip}
\usepackage{url} %to allow citing url addresses
\usepackage{rotating}
\usepackage{nameref}
\usepackage{geometry}
\usepackage[margin=1cm]{caption}
\usepackage{changepage}
\usepackage{natbib}

%define own operators&theorems&proofs
\newcommand{\expect}[1]{\mathbb{E} \left[ {#1} \right] } %define expectation symbol 
\newtheorem{proof}{Proof} %define proofs within the shadethm environment
\newcommand{\vect}[1]{\boldsymbol{#1}} %define  BOLD vectors.
\newcommand{\expe}[1]{\mathbb{E}\left[#1\right]} %define  BOLD vectors.
\newcommand{\HRule}{\rule{\linewidth}{0.01mm}} %define horizontal thin line
\newcommand{\red}[1]{{\color{red}{#1}}}

%set page margins
%\addtolength{\textheight}{2cm} 
%\addtolength{\voffset}{-1cm}
%\addtolength{\textwidth}{2cm}
%\addtolength{\hoffset}{-3cm}
\setlength{\oddsidemargin}{0pt} %set left & right margins
\setlength{\evensidemargin}{0pt} %set left & right margins
%set headers and footers
%\usepackage{fancyhdr}  
%\setlength{\headheight}{0pt}
%\pagestyle{fancy}
%\fancyhf{}
\usepackage{indentfirst}
\numberwithin{equation}{section}
%set fonts used
%\usepackage{lmodern}
%\usepackage{kpfonts}
%can check http://www.tug.dk/FontCatalogue/mathfonts.html for font alternatives
%set authoring details
\author{ %name under title %email in footnote \thanks{a.lalu@tinbergen.nl} 
	    }
\title{\vspace*{-55pt} Risk Premium Inference from (Noisy) Option Price Panels \vspace*{-35pt}}
\date{\item First version: June 16, 2018 \item Previous Version : February 21, 2024 \item This version: \today}
%start file content
\begin{document}
\maketitle
%\abstract{
%\tableofcontents
%}

\onehalfspacing

\section*{Brief Description of Research Idea}
Compare the finite sample performance of implied state vs. particle filtering when conducting joint $\mathbb{P}$-$\mathbb{Q}$	estimation on noisy option price panels focusing on the implications this has for recovered (equity) risk premium components and their error distributions/confidence bands. 

\section*{Tentative Outline of Building Blocks}
\begin{enumerate}
\item Implement an efficient simulation scheme for a Heston \citep{heston} state space simulation and associated option price panels.
\item Implement IS estimation \`{a} la \cite{boswijk2015asset}.
\item Implement particle filtering \`{a} la \cite{hurn2015pf}.
\item No noise/best case.
\item Noisy option panel case with an assorted menu of assumptions regarding option price error distributions.
\end{enumerate}

\section*{Rationale}
\begin{enumerate}
\item[A.] Noisy option data is something I had to deal with both for the first two papers, outside academia and having spent some more time looking at post 2018 academic work it seems that this is an emerging topic with a few papers that touch on this in a different context: \cite{andersen2021spatial}, \cite{duarte2024very}.
\item[B.] Topic itself ties in well with the other two papers in my thesis.
\item[C.] I have sufficient background to work on this independently.
\end{enumerate}

\newpage
 
\section{Assumed Data Generating Proces And Option Pricing}
We consider a filtered probability space $\left(\Omega, \mathcal{F}, \lbrace\mathcal{F}\rbrace_{t\geq0}, \mathbb{P} \right)$ satisfying the usual conditions\footnote{See e.g., \citet{protterstochastic}, p. 3.}. Under the historical probability measure, $\mathbb{P}$, the dynamics of the log-price process $\log(S_t)$ and of the stochastic volatility process $V_t$, which make up the state vector $X_t = [\log(S_t), V_t]$, are assumed to be described by the pair of s.d.e.:
\begin{align}
	\dif{\log(S_t)} &= \left((r_t-q_t)+ \left(\eta-1/2\right) V_t \right)\dif t + \sqrt{V_t} \dif W_{t,1}^\mathbb{P}; \label{PHestonS}\\
	\dif{V_t} &= k(\bar{V}-V_t)\dif t + \sigma_v\sqrt{V_t}\left(\rho \dif W_{t,1}^\mathbb{P} + \sqrt{1-\rho^2}\dif W_{t,2}^\mathbb{P}\right). \label{PHestonV}
\end{align}
$\left(\dif W_{t,1}^\mathbb{P}, \dif W_{t,2}^\mathbb{P}\right)$ are $\mathcal{F}_t$-adapted Brownian motion processes. $\lbrace r \rbrace_{t\geq 0}$  and  $\lbrace q \rbrace_{t\geq 0}$ are $\mathcal{F}_t$-adapted processes (e.g., deterministic) which describe the risk-free rate of interest and the continuous dividend yield respectively. $\eta$ is a parameter which governs the size of the equity risk premium. The local stochastic volatility process mean reverts at a mean reversion rate $\kappa$ to the long run mean level $\bar{V}$. $\sigma_v$ governs the volatility of volatility and $\rho$ represents the instantaneous correlation between returns and volatility.

Under the risk neutral probability measure, $\mathbb{Q}$, the model dynamics are assumed to be: 
\begin{align}
	\dif{\log(S_t)} &= \left( (r_t-q_t) - 1/2 V_t\right) \dif t + \sqrt{V_t} \dif W_{t,1}^\mathbb{Q};\label{QHestonS}\\
	\dif{V_t} &= k(\bar{V}-V_t)\dif t + \eta_{V}V_t\dif t + \sigma_v\sqrt{V_t}\left(\rho \dif W_{t,1}^\mathbb{Q} + \sqrt{1-\rho^2}\dif W_{t,2}^\mathbb{Q}\right).\label{QHestonV}
\end{align}
$\left(\dif W_{t,1}^\mathbb{Q}, \dif W_{t,2}^\mathbb{Q}\right)$ are Brownian motions under the (equivalent) martingale pricing measure $\mathbb{Q}$.

Jointly, \ref{QHestonS} and \ref{QHestonV} are known as the Heston model, which features local stochastic volatility with a leverage effect while providing semi-closed form solutions for option prices. This model is a "staple" stochastic volatility model for equity index modeling. It is used in a large number of empirical studies, has a number of well documented limitations, some of which have been remedied by enriching the model specification through the addition of jumps and/or other (latent) stochastic states. 

Given the focus of the study, attention is restricted to the simplistic set-up of the Heston model for a host of reasons. On the one hand, for most Heston model extensions, econometric approaches to parameter estimation which employ option price data come with increased computational costs. Since in the current set-up this model is known to be the data generating process, there's no potential misspecification issue at play.
Secondly, extensions come with specific estimation challenges which typically require tailored estimation designs to obtain well-behaved parameter estimates therefore limiting the comparability of results obtained throughout the study. 

\subsection{Parameter Sets}
Let $\theta$ denote the vector of parameters featured in the model. We assume $\theta \in \Theta$ where $\Theta$ is a compact set from the domain of valid data-generating process parameters. For the Heston model specification previously introduced, the value of $\theta$ corresponding to the data-generating process introduced in \eqref{PHestonS} - \eqref{QHestonV} is $\theta = \left[\eta, \kappa, \bar{v}, \sigma_v, \rho, \eta_v\right]$. Note that subsets of the parameter set play a role in either the $\mathbb{P}$-measure specification, the $\mathbb{Q}$-measure specification or in both. E.g., the Heston parameter set can be decomposed into 3 disjoint subsets: $\theta = \left[\theta^\mathbb{P}, \theta^{\mathbb{P}\cap\mathbb{Q}}, \theta^\mathbb{Q}\right]$, where $\theta^\mathbb{P} = \eta$,~ $\theta^{\mathbb{P}\cap\mathbb{Q}} = \left[\kappa, \bar{v}, \sigma_v, \rho\right]$, and $\theta^\mathbb{Q} = \eta_v$. 

\subsection{State Vector}
Assuming we have a time series of state vector observations:\footnote{In practice such a series never exists as the volatility, $V_t$, is an unobservable, i.e., latent state process} $X_1, \dots, X_T$, observed at  equally-spaced timepoints $t_1,\dots,t_T$ with $t_{i+1} - t_i = \Delta >0, \forall i\geq 1 $. Denote by $f(X_{i+1} |X_i; \theta^\mathbb{P},  \theta^{\mathbb{P} \cap \mathbb{Q}})$ the transition density of the state vector $X_t$ under the historical probability measure. Note that the transition density depends on $\theta^\mathbb{P}$ and $\theta^{\mathbb{P} \cap \mathbb{Q}}$ which represent the subset of parameters that describe the model dynamics under $\mathbb{P}$.

By using option price data the full set of parameters -- including the additional ones that characterize model dynamics under the risk neutral pricing measure -- can be estimated, i.e., inference on $\theta^\mathbb{Q}$ can be conducted while at the same time increasing sample identification potential for the $\theta^{\mathbb{P} \cap \mathbb{Q}}$ parameter set. 

\subsection{Option prices}
Suppose that, in addition to time series observations of returns, synchronous observations of option price panel data and time series data on the underlying are available. Denote the option prices by $p_{i}(T,K)$ where $T,K$ represent the maturity and the strike of the option contract. Option prices are commonly quoted in terms of their Black-Scholes implied volatility. Let $P(X_{i},T,K; \theta^{\mathbb{P}\cap\mathbb{Q}}; \theta^\mathbb{Q})$ be the corresponding model price, e.g., the option price generated by the base model specification under the equivalent martingale pricing measure for  given state vector instance $X_{i}$ and parameter set values $\theta^{\mathbb{P}\cap\mathbb{Q}}$ and $\theta^\mathbb{Q}$. If option prices were observed without error and the model is not misspecified, the observed option prices would coincide with their model-based counterparts. In practice the two prices differ. Typical modelling approaches assume away model misspecification and attribute the difference to measurement errors (e.g., microstructure noise). 

\subsubsection{Option Prices with Observation Errors}
Suppose that, at each observation date $t_i$, a panel of option prices is observed for a set of contracts
$\mathcal{C}_i=\{(T_i^{(j)},K_i^{(j)})\}_{j=1}^{q_i}$ and denote the resulting price vector by
$p_i^{\mathrm{obs}}=\left(p_i^{(1),\mathrm{obs}},\dots,p_i^{(q_i),\mathrm{obs}}\right)^\top$, where
$p_i^{(j),\mathrm{obs}}=p_i^{\mathrm{obs}}(T_i^{(j)},K_i^{(j)})$. 

Similarly, let
$p_i^{\mathrm{model}}=\left(p_i^{(1),\mathrm{model}},\dots,p_i^{(q_i),\mathrm{model}}\right)^\top$
be the vector of corresponding model-based prices implied by $P(X_i,T,K;\theta^{\mathbb{P}\cap\mathbb{Q}};\theta^\mathbb{Q})$.
The difference between $p_i^{\mathrm{obs}}$ and $p_i^{\mathrm{model}}$ is interpreted as option observation error.
This difference can be, inter alia, attributed to microstructure effects, stale quotes or price discreteness due to minimum tick sizing or quote rounding conventions. \citet{duarte2024very}
suggests that such frictions can be quantitatively important for inference involving option-implied risk premia. 
 
The proposed simulation design therefore does not commit ex-ante to a single parametric error law.
Instead, it allows the option panel observation errors to be generated from a menu of nested families, which evolve in complexity along two dimensions. First, it allows for simple benchmark perturbations which treat errors as independent across contracts (including additive or multiplicative representations) while permitting the marginal scale of the errors to vary with moneyness and maturity. This is motivated by the fact that option observation errors
can enter the asymptotic behavior of option-based estimators through contracts near the spot, and hence local
features of the panel can matter disproportionately. \citet{chongtodorov2024vov}

Second, the design allows for dependence across the option panel in the sense that errors may be correlated across nearby strikes and tenors. This is motivated by econometric evidence and formal testing procedures that explicitly address spatial dependence in option observation errors as discussed in \citet{andersen2021spatial}.

In addition, it allows for the possibility that the error structure varies over time through  common components (e.g., akin to persistent quoting regimes), which provides a parsimonious route to temporal dependence without forcing a full scale spatio-temporal observation error dependence.

The detailed distributional assumptions and the contract design $\mathcal{C}_i$ are deferred to the simulation design section, where the different error families are used to stress-test the estimator comparison under progressively more realistic option panel noise.

\subsubsection{Option Prices with Observation Errors}
Suppose that, at each observation date $t_i$, a panel of option prices is observed for a set of contracts
$\mathcal{C}_i=\{(T_i^{(j)},K_i^{(j)},\mathrm{type}^{(j)})\}_{j=1}^{q_i}$ and denote the resulting price vector by
$p_i^{\mathrm{obs}}=\left(p_i^{(1),\mathrm{obs}},\dots,p_i^{(q_i),\mathrm{obs}}\right)^\top$, where
$p_i^{(j),\mathrm{obs}}=p_i^{\mathrm{obs}}(T_i^{(j)},K_i^{(j)})$.

Similarly, let
$p_i^{0}=p_i^{\mathrm{model}}=\left(p_i^{(1),0},\dots,p_i^{(q_i),0}\right)^\top$
be the vector of corresponding model-based prices implied by
$P(X_i,T,K;\theta^{\mathbb{P}\cap\mathbb{Q}};\theta^\mathbb{Q})$.
Let $\sigma_i^{0}=\left(\sigma_i^{(1),0},\dots,\sigma_i^{(q_i),0}\right)^\top$
denote the corresponding Black-Scholes implied volatility vector and let
$\mathcal{V}_i=\left(\mathcal{V}_i^{(1)},\dots,\mathcal{V}_i^{(q_i)}\right)^\top$
denote the associated Black-Scholes vega vector.

The difference between $p_i^{\mathrm{obs}}$ and $p_i^{0}$ is interpreted as option observation error.
This difference can be, inter alia, attributed to microstructure effects, price discreteness due to minimum
tick sizing or quote rounding conventions, and errors arising from the construction of option price surfaces.
\citet{duarte2024very} suggests that such frictions can be quantitatively important for inference involving
option-implied risk premia.

The proposed simulation design therefore considers three observed-panel specifications generated from the same
clean model-implied panel. For $j=1,\ldots,q_i$, let
\[
m_i^{(j)}=\log\left(\frac{K_i^{(j)}}{F_i^{(j)}}\right),
\qquad
\tau_i^{(j)}=T_i^{(j)}-t_i,
\]
where $F_i^{(j)}$ denotes the corresponding forward price. The marginal implied-volatility error scale is
specified as
\begin{align}
s_\sigma(m_i^{(j)},\tau_i^{(j)})
=
\alpha_0\left(
1+\alpha_m |m_i^{(j)}|
+\frac{\alpha_\tau}{\sqrt{\max(\tau_i^{(j)},\tau_{\min})}}
\right),
\label{eq:iv_error_scale}
\end{align}
and the corresponding price error scale is defined by
\begin{align}
s_{p,i}^{(j)}
=
\max\left\{
\chi_{\mathrm{tick}}\delta_{\mathrm{tick}},
\max\left(\mathcal{V}_i^{(j)},\mathcal{V}_{\min}\right)
s_\sigma(m_i^{(j)},\tau_i^{(j)})
\right\}.
\label{eq:price_error_scale}
\end{align}
The first term in \eqref{eq:price_error_scale} imposes a tick-size based lower bound on the price perturbation.
The second term converts an implied-volatility perturbation into a price perturbation through the local Black-Scholes
vega.

The first specification is a liquid-market benchmark in which the option price errors are independent across
contracts conditional on the clean panel:
\begin{align}
\tilde p_i^{(j),A}
=
p_i^{(j),0}+s_{p,i}^{(j)} Z_i^{(j)},
\qquad
Z_i^{(j)}\overset{i.i.d.}{\sim}\mathcal{N}(0,1).
\label{eq:noise_liquid}
\end{align}
This specification is intended to represent small quote perturbations whose scale varies with moneyness, maturity
and tick size, without imposing dependence across contracts.

The second specification allows for cross-sectional dependence across the option panel. Let
$\varepsilon_i^B=\left(\varepsilon_i^{(1),B},\ldots,\varepsilon_i^{(q_i),B}\right)^\top$ and set
\begin{align}
\tilde p_i^{B}=p_i^{0}+\varepsilon_i^B,
\qquad
\varepsilon_i^B\sim\mathcal{N}(0,D_i R_i D_i),
\label{eq:noise_spatial}
\end{align}
where
\[
D_i=\mathrm{diag}\left(s_{p,i}^{(1)},\ldots,s_{p,i}^{(q_i)}\right)
\]
and the correlation matrix $R_i$ is given by
\begin{align}
R_{i,jk}
=
\exp\left(
-\frac{|m_i^{(j)}-m_i^{(k)}|}{\ell_m}
-\frac{|\tau_i^{(j)}-\tau_i^{(k)}|}{\ell_\tau}
\right).
\label{eq:spatial_kernel}
\end{align}
This specification captures the possibility that nearby strikes and maturities share common quote or surface
construction errors. This is consistent with the econometric evidence on spatial dependence in option observation
errors documented in \citet{andersen2021spatial}.

The third specification introduces low-dimensional persistent surface distortions. Let
$f_i=(f_{0,i},f_{1,i},f_{2,i})^\top$ evolve according to
\begin{align}
f_i=A f_{i-1}+\xi_i,
\qquad
\xi_i\sim\mathcal{N}(0,Q),
\label{eq:factor_noise}
\end{align}
where $A$ is diagonal with entries inside the unit circle. For contract $j$, define
\[
b_i^{(j)}=\left(1,m_i^{(j)},\tau_i^{(j)}\right)^\top.
\]
The implied-volatility distortion is
\begin{align}
\Delta\sigma_i^{(j),C}
=
\left(b_i^{(j)}\right)^\top f_i + u_i^{(j)},
\qquad
u_i^{(j)}\sim\mathcal{N}\left(0,\left(s_{\sigma,u}(m_i^{(j)},\tau_i^{(j)})\right)^2\right),
\label{eq:factor_iv_noise}
\end{align}
and the corresponding raw price perturbation is
\begin{align}
\tilde p_i^{(j),C}
=
p_i^{(j),0}
+
\max\left(\mathcal{V}_i^{(j)},\mathcal{V}_{\min}\right)
\Delta\sigma_i^{(j),C}.
\label{eq:factor_price_noise}
\end{align}
The factors in \eqref{eq:factor_noise} are intended to represent persistent level, moneyness-slope and maturity-slope
distortions in the quoted surface.

For all three specifications, the raw contaminated price $\tilde p_i^{(j)}$ is rounded to the relevant tick size
and projected back into the no-arbitrage price interval for the corresponding European option. The observed implied
volatility $\sigma_i^{(j),\mathrm{obs}}$ is then obtained by inverting the Black-Scholes pricing formula at the
resulting observed price. This final step ensures that the contaminated panel contains both observed prices and
observed implied volatilities, and that invalid prices generated by the noise mechanism are handled by capping rather
than by re-sampling.

\emph{Example sampling noise for specification A: [to be added].}

\emph{Example sampling noise for specification B: [to be added].}

\emph{Example sampling noise for specification C: [to be added].}

\subsubsection{Particle Filter Based Estimation}
Estimation contexts which make use of option prices (that are nonlinear, often semi-analytic expressions featuring model parameters and partially observable state vectors) can be computationally demanding. The models used to describe index returns for option pricing purposes typically feature latent stochastic state processes which require filtering  or implying. This increases the computational cost of econometric estimation exercises as one has to devise parsimonious estimators which can work with non-analytic expressions of state vector observations by relying on stable and robust numerical procedures. 

%The advent of parallel computing techniques has rendered advanced particle filter methods feasible for set-ups with larger panels of option price observations [references, some of them very new indeed], allowing for the evaluation of (quasi-)likelihood functions of option price observations.

Two other hurdles are typically encountered when working with continuous time models. Firstly, even the simple continuous time stochastic model specifications (like the base-case model used in this paper) do not have known densities readily usable to evaluate the likelihood function. Secondly, one often has to resort to using discretisation schemes to simulate state vector (i.e., particle) paths. However, despite this second hurdle, approximating transition densities using discrete distributions consisting of particle support points provides a tractable workaround which can easily be adapted for most of the typical model specifications considered. The filtering distributions provide estimators for the latent states and their possible distributions. The Monte Carlo procedures involved in the simulation of particles can be tailored to accommodate models with jumps and with more complex state dynamics than featured in the Heston model specification. 

In the context of the base model specification previously introduced, we outline the following optimal filtering procedure in the spirit of \citet{johannes2009optimal}, \citet{christoffersen2010volatility} and ***[Hurn et al]*** and introduce the necessary implementation details for the Monte Carlo estimation exercise. First, let us denote the  log-likelihood function of the model under the physical probability measure by
\begin{equation}
\log \mathcal{L}_T = \log f(X_0) + \sum_{i=0}^T \log f(X_{i+1}|X_{i};\theta^{\mathbb{P},\mathbb{P}\cap\mathbb{Q}}).
\end{equation}
The evaluation of this likelihood function is unfeasible as the state vector is not fully observed. Also note that it does not provide any information about $\theta^\mathbb{Q}$, the parameters which feature in the dynamics of the model under the risk neutral probability measure.
 
To take into account option prices and to allow for the likelihood to be evaluated, a filtering rule using the time series observations for the components of the state vector which are observable is introduced alongside an option price panel likelihood observation in order to extract information about the latent states which are not observable.  Let us denote by $\mathcal{O}_i$ the filtration containing the collection of sigma algebras generated by the time series of observable states and associated option prices. E.g., for our base model specification this is the filtration generated by $S_{1},\dots,S_{T}$ and the panel $p_{1}(T,K),\dots  p_{T}(T,K)$ of option prices for all available option maturities $T$ and strike prices $K$. The level of the stochastic volatility process $V_{i}$ is unobservable. To evaluate the likelihood contribution, the following need to be evaluated:\vspace{+5pt}\\
\texttt{Step 1. $T_{i+1}$ likelihood contribution:} \vspace{-5pt}
\begin{align}
f(X_{i+1}|\mathcal{O}_{i};\theta) &=\int_{\mathcal{D}_{V}}  f(S_{i+1}|{V}_{i+1},\mathcal{O}_i;\theta) f(V_{i+1}|\mathcal{O}_i;\theta) \dif{V_{i+1}} \notag\\
&=\int_{\mathcal{D}_{V}} f(S_{i+1}|{V}_{i+1},\mathcal{O}_i; \theta) \int_{\mathcal{D}_{V}} f(V_{i+1}|{V}_{i},\mathcal{O}_i;\theta) \notag \\
& \hspace{165pt} \times f(V_{i}|\mathcal{O}_i; \theta)  \dif V_{i} \dif{V_{i+1}}\label{LikelihoodContribTheo}
\end{align}

\texttt{Step 2. Latent state variable density update:}\vspace{-5pt}
\begin{align}
f(V_{i+1}|\mathcal{O}_{i+1}; \theta) &= \frac{f(X_{i+1},V_{i+1}|\mathcal{O}_{i+1}; \theta)}{f(X_{i+1}|\mathcal{O}_{i+1}; \theta)} \label{SamplingDensityTheo}
\end{align}
%
%To evaluate \ref{LikelihoodContribTheo} and \ref{SamplingDensityTheo} numerically a combination of Monte Carlo path draws and Monte Carlo integrations is used. Starting with an initial cloud of particles (i.e., a set of likely values for ${V_i}$ pre-sampled from $\mathcal{D}_{V}$) the state vector $X_{i} = \left[S_i, V_i\right]$ is advanced using a discretisation approach (e.g., an Euler scheme). The next step requires calculating $•$

\newpage
%========================================================
\appendix
\section{Simulation of Heston Paths and Model Option Panels}
\label{app:sim_heston}

This appendix documents the simulation procedure used to generate paths of the state vector
$X_t=[\log(S_t),V_t]$ under the historical probability measure $\mathbb{P}$ and the associated
panel of model-implied European option prices under the equivalent martingale measure $\mathbb{Q}$,
based on the data-generating process \eqref{PHestonS}--\eqref{PHestonV} and pricing dynamics
\eqref{QHestonS}--\eqref{QHestonV}.

\subsection{Sampling grid and contract set}
Let $\{t_i\}_{i=0}^T$ denote an equally-spaced time grid with $t_{i+1}-t_i=\Delta>0$.
The simulated sample consists of $\{(\log(S_{t_i}),V_{t_i})\}_{i=0}^T$.
At each date $t_i$, we define a (possibly time-varying) set of option contracts
$\mathcal{C}_i=\{(T_i^{(j)},K_i^{(j)},\mathrm{type}^{(j)})\}_{j=1}^{q_i}$ and compute a vector of
model-implied prices $p_i^{\mathrm{model}}\in\mathbb{R}^{q_i}$.

\subsection{State simulation under $\mathbb{P}$}
Let $y_t:=\log(S_t)$ and write $y_i:=y_{t_i}$ and $V_i:=V_{t_i}$. Under $\mathbb{P}$,
the dynamics are given by \eqref{PHestonS}--\eqref{PHestonV}.
We simulate $V_{i+1}$ using the exact transition of the square-root process and then update $y_{i+1}$
using a discretisation that preserves the leverage channel implied by $\rho$.

\subsubsection{Exact transition for $V_{i+1}\mid V_i$}
Conditional on $V_i$, the square-root process in \eqref{PHestonV} admits an exact transition
representation in terms of a scaled noncentral chi-square random variable:
\begin{align}
V_{i+1} \;\overset{d}{=}\; c\,\chi'^2_{\nu}(\lambda),
\label{eq:CIR_exact}
\end{align}
where
\begin{align}
c &= \frac{\sigma_v^2\left(1-e^{-\kappa \Delta}\right)}{4\kappa}, &
\nu &= \frac{4\kappa \bar{V}}{\sigma_v^2}, &
\lambda &= \frac{4\kappa e^{-\kappa \Delta}}{\sigma_v^2\left(1-e^{-\kappa \Delta}\right)}V_i.
\label{eq:CIR_exact_params}
\end{align}
This step guarantees $V_{i+1}\geq 0$ and removes Euler-type discretization bias in the variance
marginal transition.

\subsubsection{Integrated variance proxy}
Define the integrated variance over $[t_i,t_{i+1}]$ by
\begin{align}
I_i := \int_{t_i}^{t_{i+1}} V_s\,\dif s.
\end{align}
Since $I_i$ is not sampled explicitly, we use the trapezoidal proxy
\begin{align}
\widehat{I}_i := \frac{\Delta}{2}\left(V_i+V_{i+1}\right).
\label{eq:I_trap}
\end{align}

\subsubsection{Log-price update with leverage}
Let $Z_{i+1}\sim\mathcal{N}(0,1)$ i.i.d. across $i$ and independent of the variance draw.
Using the decomposition of the Brownian drivers in \eqref{PHestonS}--\eqref{PHestonV},
the log-price increment is simulated as
\begin{align}
\Delta y_i
&:= y_{i+1}-y_i \nonumber\\
&= (r_{t_i}-q_{t_i})\Delta
+ \left(\eta-\frac{1}{2}\right)\widehat{I}_i
+ \frac{\rho}{\sigma_v}\left[\left(V_{i+1}-V_i\right)-\kappa\left(\bar{V}\Delta-\widehat{I}_i\right)\right]
+ \sqrt{1-\rho^2}\sqrt{\widehat{I}_i}\,Z_{i+1}.
\label{eq:logS_update}
\end{align}
Finally, set $y_{i+1}=y_i+\Delta y_i$ and $S_{i+1}=\exp(y_{i+1})$.

\subsection{Model option panels under $\mathbb{Q}$}
Given the simulated state $X_i=[\log(S_i),V_i]$ at time $t_i$, model option prices are computed under
$\mathbb{Q}$ for each contract $(T_i^{(j)},K_i^{(j)})\in\mathcal{C}_i$ by
\begin{align}
p_{i}^{\mathrm{model}}(T_i^{(j)},K_i^{(j)})
= P\!\left(X_i,T_i^{(j)},K_i^{(j)};\theta^{\mathbb{P}\cap\mathbb{Q}},\theta^{\mathbb{Q}}\right),
\qquad j=1,\dots,q_i,
\label{eq:panel_model_prices}
\end{align}
with $\theta=[\eta,\kappa,\bar V,\sigma_v,\rho,\eta_V]$ and the partitioning used in the main text.
Under $\mathbb{Q}$, the variance drift \eqref{QHestonV} can be written in standard Heston form by defining
\begin{align}
\kappa_{\mathbb{Q}} := \kappa-\eta_V,
\qquad
\bar V_{\mathbb{Q}} := \frac{\kappa\bar V}{\kappa-\eta_V},
\end{align}
so that $dV_t=\kappa_{\mathbb{Q}}(\bar V_{\mathbb{Q}}-V_t)\,dt+\sigma_v\sqrt{V_t}\,dW^{\mathbb{Q}}_{t,v}$.

\subsubsection{COS pricing}
European option prices in \eqref{eq:panel_model_prices} are evaluated using the Fourier-cosine (COS)
method, which expands the risk-neutral density of the log-price (or log-return) on a bounded interval
as a cosine series and integrates the payoff against this approximation.

Let $\tau:=T-t_i$ and define the log-forward moneyness
\begin{align}
x:=\log\!\left(\frac{S_i e^{(r_{t_i}-q_{t_i})\tau}}{K}\right).
\end{align}
Denote by $\phi(u;\tau,X_i)$ the conditional characteristic function of $\log(S_T)$ under $\mathbb{Q}$
(which is available in semi-closed form for Heston). The COS approximation for a European call can be
written as
\begin{align}
C_i(T,K) \approx e^{-r_{t_i}\tau}\sum_{k=0}^{N-1}{}'
\Re\!\left\{\phi\!\left(\frac{k\pi}{b-a};\tau,X_i\right)\exp\!\left(-i\frac{k\pi a}{b-a}\right)\right\}
\,U_k(x),
\label{eq:COS_call}
\end{align}
where $[a,b]$ is a truncation range for the log-price (or log-return) support, $N$ is the number of cosine
terms, and $\sum'$ indicates that the $k=0$ term is weighted by $1/2$. The payoff coefficients $U_k(x)$
admit closed forms for vanilla calls and puts. Put prices are obtained either via the corresponding COS payoff
coefficients or via put-call parity.

\section{Option Panel Observation Error Specifications}
\label{app:option_noise_specs}

This appendix documents the functional forms used to generate contaminated option panels from the clean model-implied
option panel. The purpose is to introduce observation error structures which are simple enough to be implemented in a
large-scale Monte Carlo design, but sufficiently close to observed option-market features to provide a meaningful
stress-test of PF and IS estimation.

\subsection{Marginal error scale}
The marginal error scale is specified in implied-volatility units by
\begin{align}
s_\sigma(m,\tau)
=
\alpha_0\left(
1+\alpha_m |m|
+\frac{\alpha_\tau}{\sqrt{\max(\tau,\tau_{\min})}}
\right).
\label{eq:app_iv_error_scale}
\end{align}
The term $\alpha_0$ controls the baseline magnitude of quote perturbations. The term $|m|$ measures the distance
from the centre of the option surface. Its inclusion reflects the fact that wing options are typically less liquid,
have lower vega and are more sensitive to small price perturbations when represented in implied-volatility units.
The term $\tau^{-1/2}$ is included because a fixed price perturbation maps into a larger implied-volatility
perturbation for short-dated options. In particular, for small perturbations,
\[
\Delta\sigma \approx \frac{\Delta p}{\mathcal{V}},
\]
and the Black-Scholes vega of an at-the-money option is locally proportional to $\sqrt{\tau}$. Hence, a fixed
price error corresponds approximately to an implied-volatility error of order $\tau^{-1/2}$.

The price-level scale used for price contamination is therefore
\begin{align}
s_p
=
\max\left\{
\chi_{\mathrm{tick}}\delta_{\mathrm{tick}},
\max(\mathcal{V},\mathcal{V}_{\min})s_\sigma(m,\tau)
\right\}.
\label{eq:app_price_error_scale}
\end{align}
This specification keeps the error in price units, which is the object that is contaminated in the simulation
design, while preserving the interpretation of the perturbation in implied-volatility space. The tick-size floor
ensures that the error scale does not vanish for low-vega contracts and incorporates the effect of price discreteness
induced by minimum tick sizes and quote rounding conventions. Bid-ask spreads and tick sizes are standard
microstructure mechanisms through which transaction costs and price discreteness enter traded prices
\citep{demsetz1968cost}.

\subsection{Independent liquid-market quote noise}
The first specification is
\begin{align}
\tilde p_i^{(j),A}
=
p_i^{(j),0}+s_{p,i}^{(j)}Z_i^{(j)},
\qquad
Z_i^{(j)}\overset{i.i.d.}{\sim}\mathcal{N}(0,1).
\label{eq:app_liquid_noise}
\end{align}
This produces independent price perturbations across contracts, while allowing the marginal scale to vary with
moneyness, maturity, vega and tick size. This specification serves as the low-noise benchmark. It preserves the
idea that the clean panel is nearly observed, while avoiding the unrealistic implication that all contracts are
observed with the same precision.

\emph{Example sampling noise: [to be added after calibration of $(\alpha_0,\alpha_m,\alpha_\tau,\chi_{\mathrm{tick}})$].}

\subsection{Spatially correlated surface noise}
The second specification is
\begin{align}
\tilde p_i^B=p_i^0+\varepsilon_i^B,
\qquad
\varepsilon_i^B\sim\mathcal{N}(0,D_iR_iD_i),
\label{eq:app_spatial_noise}
\end{align}
where
\[
D_i=\mathrm{diag}\left(s_{p,i}^{(1)},\ldots,s_{p,i}^{(q_i)}\right)
\]
and
\begin{align}
R_{i,jk}
=
\exp\left(
-\frac{|m_i^{(j)}-m_i^{(k)}|}{\ell_m}
-\frac{|\tau_i^{(j)}-\tau_i^{(k)}|}{\ell_\tau}
\right).
\label{eq:app_spatial_kernel}
\end{align}
This functional form imposes a local dependence structure on the option surface. Two contracts that are close in
log-forward moneyness and maturity have highly correlated observation errors, whereas contracts far apart on the
surface have weakly correlated observation errors. The parameters $\ell_m$ and $\ell_\tau$ determine the rate at
which the correlation decays across moneyness and maturity respectively. The exponential form is used because it
allows local dependence to decay monotonically with distance without imposing an overly smooth error field.

This specification is directly motivated by the evidence in \citet{andersen2021spatial}, which develops tests for
spatial dependence in option observation errors. It also reflects the practical fact that observed option panels are
typically produced as surfaces rather than as collections of unrelated contracts. Nearby strikes and maturities may
therefore share interpolation, smoothing, quoting or broker-source errors.

\emph{Example sampling noise: [to be added after calibration of $(\ell_m,\ell_\tau)$ and the marginal scale parameters].}

\subsection{Persistent low-rank surface distortions}
The third specification is built from persistent low-dimensional surface factors. Let
\begin{align}
f_i=A f_{i-1}+\xi_i,
\qquad
\xi_i\sim\mathcal{N}(0,Q),
\label{eq:app_surface_factor}
\end{align}
where $f_i=(f_{0,i},f_{1,i},f_{2,i})^\top$ and $A$ is diagonal. For contract $j$, define
\[
b_i^{(j)}=(1,m_i^{(j)},\tau_i^{(j)})^\top.
\]
The implied-volatility distortion is
\begin{align}
\Delta\sigma_i^{(j),C}
=
\left(b_i^{(j)}\right)^\top f_i + u_i^{(j)},
\qquad
u_i^{(j)}\sim\mathcal{N}\left(0,\left(s_{\sigma,u}(m_i^{(j)},\tau_i^{(j)})\right)^2\right),
\label{eq:app_factor_iv_noise}
\end{align}
and the price-level perturbation is
\begin{align}
\tilde p_i^{(j),C}
=
p_i^{(j),0}
+
\max\left(\mathcal{V}_i^{(j)},\mathcal{V}_{\min}\right)
\Delta\sigma_i^{(j),C}.
\label{eq:app_factor_price_noise}
\end{align}
The first component of $f_i$ acts as a level distortion of the implied-volatility surface. The second component
acts as a moneyness-slope distortion and is therefore related to skew. The third component acts as a maturity-slope
distortion. This factor specification is consistent with the empirical and practical treatment of implied volatility
surfaces as low-dimensional objects whose dominant movements can be summarized by level, slope and curvature-like
directions \citep{cont2002dynamics}. The persistence in \eqref{eq:app_surface_factor} captures the fact that surface
construction errors and quote-quality conditions need not be independent across adjacent observation dates.

\emph{Example sampling noise: [to be added after calibration of $(A,Q)$ and the residual scale parameters].}

\subsection{Tick rounding, price validity and implied-volatility inversion}
For each raw contaminated price $\tilde p_i^{(j)}$, tick rounding is applied through
\[
\mathcal{R}_{\delta_{\mathrm{tick}}}(x)
=
\delta_{\mathrm{tick}}
\mathrm{round}\left(\frac{x}{\delta_{\mathrm{tick}}}\right).
\]
Let $L_i^{(j)}$ and $U_i^{(j)}$ denote the no-arbitrage lower and upper bounds for the corresponding European
option. For a call option,
\[
L_i^{(j)}=\max\left(S_ie^{-q\tau_i^{(j)}}-K_i^{(j)}e^{-r\tau_i^{(j)}},0\right),
\qquad
U_i^{(j)}=S_ie^{-q\tau_i^{(j)}}.
\]
For a put option,
\[
L_i^{(j)}=\max\left(K_i^{(j)}e^{-r\tau_i^{(j)}}-S_ie^{-q\tau_i^{(j)}},0\right),
\qquad
U_i^{(j)}=K_i^{(j)}e^{-r\tau_i^{(j)}}.
\]
The observed price is obtained by projecting the rounded price into the admissible interval:
\begin{align}
p_i^{(j),\mathrm{obs}}
=
\min\left\{
U_i^{(j)}-\epsilon_p,
\max\left[
L_i^{(j)}+\epsilon_p,
\mathcal{R}_{\delta_{\mathrm{tick}}}\left(\tilde p_i^{(j)}\right)
\right]
\right\}.
\label{eq:app_price_projection}
\end{align}
The projection in \eqref{eq:app_price_projection} is used instead of re-sampling. This keeps the distributional
draws reproducible and preserves information about the frequency with which the noise mechanism generates boundary
prices. The implementation therefore records whether capping is applied at the lower or upper bound.

The observed Black-Scholes implied volatility $\sigma_i^{(j),\mathrm{obs}}$ is computed from
$p_i^{(j),\mathrm{obs}}$ using a robust implied-volatility inversion routine. In the implementation this is handled
through a stand-alone implied-volatility solver interface so that the same routine can be used for clean panels,
contaminated panels and later estimator diagnostics. The preferred numerical method is the rational approximation
approach of \citet{jaeckel2015lets}, which is designed for fast and robust implied-volatility inversion.

\subsection{Implementation of contaminated panels}
The clean model-implied panels are generated once and stored before any observation error is applied. Each
contaminated panel is then generated as a deterministic transformation of the clean panel, conditional on the
noise specification and the random seed assigned to the sample and scenario. This avoids re-pricing the Heston
model for each observation-error design.

The implementation records, for each contract, the clean model price, the clean Black-Scholes implied volatility,
the raw contaminated price before validity enforcement, the rounded and capped observed price, the observed implied
volatility, the noise scenario and indicators for whether lower- or upper-bound capping was applied. The capping
policy is preferred to re-sampling because it preserves reproducibility and makes the frequency of invalid raw
noise draws observable.

The same implied-volatility inversion routine is used for clean and contaminated panels. This is important because
the subsequent IS and PF estimation exercises may use observed implied volatilities, price-space errors, or
vega-scaled price errors. A single stand-alone implied-volatility solver therefore prevents the construction of
different data sets from depending on different numerical inversion conventions.

%========================================================
\newpage

\bibliographystyle{apalike}
\nocite{*}
\bibliography{P_ONE_BIB}
\end{document}